\SixthArticle{48}

It will be worthwhile to show how, with the aid of transformations 82,91, all the formulas collected in Article 5 can be deduced very easily from the binomial theorem alone. Formulas I---IV follow from it immediately. Formulas VI---IX follow of themselves if \(e^x\) is regarded as the limit of the power \(\left(1+\frac{x}{i}\right)^i\) or \(\left(1-\frac{x}{i}\right)^{-i}\), and \(\log x\) as the limit of \(i(x^{\frac{1}{i}}-1)\), as \(i\) increases to infinity. Further, from
\[
\begin{aligned}[c]
\cos n\varphi+\BareRoot{-1}.\sin n\varphi&=(\cos\varphi+\BareRoot{-1}.\sin\varphi)^n\\
\cos n\varphi-\BareRoot{-1}.\sin n\varphi&=(\cos\varphi-\BareRoot{-1}.\sin\varphi)^n
\end{aligned}
\]
formulas XVIII and XXII follow by subtraction and addition, and from these, by formula 82, XIX and XXIII immediately; from these again, by formula 91, XVI, XVII, XX and XXI are deduced. Setting \(t\) in place of \(nt\) and taking \(n\) infinite gives XI and XII from XVI and XX; taking \(n\) infinitely small gives XIII---XV from XVI---XVIII.

\SixthArticle{49}

The substitution \(x=\frac{1}{y}\) similarly gives
\[
\begin{aligned}&0=\alpha\beta P-(\alpha+\beta-1-(\gamma-2)y)y\frac{\mathrm{d}P}{\mathrm{d}y}\\&+(yy-y^3)\frac{\mathrm{d}\mathrm{d}P}{\mathrm{d}y^2}\end{aligned}
\]
Then setting \(P=y^\mu P'\) gives
\[
\begin{aligned}[c]
&\mathrm{I.}\qquad 0=P'(\alpha\beta-\mu(\alpha+\beta-1)\\&\qquad+\mu(\gamma-2)y+(\mu\mu-\mu)(1-y))\\
&-\frac{\mathrm{d}P'}{\mathrm{d}y}(\alpha+\beta-1-(\gamma-2)y-2\mu(1-y))y\\
&+(yy-y^3)\frac{\mathrm{d}\mathrm{d}P'}{\mathrm{d}y^2}
\end{aligned}
\]
For the multiplier of \(P'\) to become divisible by \(y\), either \(\mu=\alpha\) or \(\mu=\beta\) must be set; the former value produces
\[
\begin{aligned}&\mathrm{II.}\qquad
0=P'\alpha(\gamma-\alpha-1)-\frac{\mathrm{d}P'}{\mathrm{d}y}(\beta-\alpha-1-(\gamma-2\alpha-2)y)\\&+(y-yy)\frac{\mathrm{d}\mathrm{d}P'}{\mathrm{d}y^2}\end{aligned}
\]
whose \textit{particular} solution is
\[
P'=F(\alpha,\alpha+1-\gamma,\alpha+1-\beta,y)
\]
\par\SourceLeaf{229}
Thus equation I is satisfied by the particular solution
\[
P=y^\alpha F(\alpha,\alpha+1-\gamma,\alpha+1-\beta,y).
\]
and in the same way the other value \(\mu=\beta\) supplies another particular solution
\[
P=y^\beta F(\beta,\beta+1-\gamma,\beta+1-\alpha,y)
\]
from which the general solution is obtained:
\[
\begin{aligned}&P=Ay^\alpha F(\alpha,\alpha+1-\gamma,\alpha+1-\beta,y)\\&+By^\beta F(\beta,\beta+1-\gamma,\beta+1-\alpha,y)\end{aligned}
\]
\Needspace{11\baselineskip}
where \(A,B\) denote constants, which, however, are not arbitrary but fully determined, since \(P\) is not the general solution of equation 80, but only a particular one. But so that determining the values of the constants \(A,B\) does not lead us into useless detours, we shall deduce the same equation in another way from what we have already developed.

In equation 91, setting \(-\frac{x}{1-x}=1-z\), and replacing in equation 86 \(\beta\) by \(\gamma-\beta\), \(\gamma\) by \(\alpha+1-\beta\), and \(x\) by \(z\), gives
\[
\begin{aligned}[c]
&(1-x)^\alpha F(\alpha,\beta,\gamma,x)\\&=
\frac{\Pi(\gamma-1)\Pi(\beta-\alpha-1)}{\Pi(\gamma-\alpha-1)\Pi(\beta-1)}F(\alpha,\gamma-\beta,\alpha+1-\beta,z)\\
&+\frac{\Pi(\gamma-1)\Pi(\alpha-\beta-1)}{\Pi(\alpha-1)\Pi(\gamma-\beta-1)}z^{\beta-\alpha}F(\beta,\gamma-\alpha,\beta+1-\alpha,z)
\end{aligned}
\]
But by equations 91, 92,
\[
\begin{aligned}[c]
&F(\alpha,\gamma-\beta,\alpha+1-\beta,z)\\&=(1-z)^{-\alpha}F\left(\alpha,\alpha+1-\gamma,\alpha+1-\beta,-\frac{z}{1-z}\right)\\[5pt]
&F(\beta,\gamma-\alpha,\beta+1-\alpha,z)\\&=(1-z)^{-\beta}F\left(\beta,\beta+1-\gamma,\beta+1-\alpha,-\frac{z}{1-z}\right)
\end{aligned}
\]
On substituting these, as well as \(z=\frac{1}{1-x}\), \(1-z=-\frac{x}{1-x}\), \(-\frac{z}{1-z}=\frac{1}{x}\), we have
\[
\begin{aligned}[c]
&[93]\qquad F(\alpha,\beta,\gamma,x)\\&=
\frac{\Pi(\gamma-1)\Pi(\beta-\alpha-1)}{\Pi(\gamma-\alpha-1)\Pi(\beta-1)}(-x)^{-\alpha}F\left(\alpha,\alpha+1-\gamma,\alpha+1-\beta,\frac{1}{x}\right)\\
&+\frac{\Pi(\gamma-1)\Pi(\alpha-\beta-1)}{\Pi(\alpha-1)\Pi(\gamma-\beta-1)}(-x)^{-\beta}F\left(\beta,\beta+1-\gamma,\beta+1-\alpha,\frac{1}{x}\right)
\end{aligned}
\]
which agrees with the equation found above on setting
\[
\begin{aligned}[c]
A&=\frac{\Pi(\gamma-1)\Pi(\beta-\alpha-1)}{\Pi(\gamma-\alpha-1)\Pi(\beta-1)}(-1)^\alpha\\
B&=\frac{\Pi(\gamma-1)\Pi(\alpha-\beta-1)}{\Pi(\alpha-1)\Pi(\gamma-\beta-1)}(-1)^\beta
\end{aligned}
\]
\SourceLeaf{230}
where it should be noted that
\[
\begin{array}{rcl}
(-1)^\alpha&=&\cos\alpha k\pi+\BareRoot{-1}\,.\sin\alpha k\pi\\
(-1)^\beta&=&\cos\beta k\pi+\BareRoot{-1}\,.\sin\beta k\pi
\end{array}
\]
with \(k\) denoting any odd integer without further specification.

\SixthArticle{50}
By equation 93, the value of our function for values of the fourth element greater than unity is reduced to the case where the fourth element is less than unity. At the same time, it is evident that \textit{negative} values of the fourth element greater than unity always correspond to one real value of the function \(F\), whereas positive values can correspond to a real value of the function only when \(\alpha\) and \(\beta\) are either integers or rational fractions whose denominators are odd; in the remaining cases, \(F(\alpha,\beta,\gamma,x)\) admits only imaginary values for a positive value of \(x\) greater than unity.

\SixthArticle{51}
The relations among several functions \(F\) developed so far have all been linear: we add another of a different kind. Let
\[
\begin{array}{rcl}
P&=&F(\alpha,\beta,\gamma,x)\\
Q&=&x^{1-\gamma}F(\alpha+1-\gamma,\beta+1-\gamma,2-\gamma,x)\\
R&=&F(\alpha,\beta,\alpha+\beta+1-\gamma,1-x)
\end{array}
\]
so that \(P,\ Q.\ R\) are three particular solutions of equation 80, or
\[
\begin{aligned}
\mathrm{I.}\quad&0=\alpha\beta P-(\gamma-(\alpha+\beta+1)x)\dfrac{\mathrm dP}{\mathrm dx}\\&\qquad-(x-xx)\dfrac{\mathrm{dd}P}{\mathrm dx^2}\\[4pt]
\mathrm{II.}\quad&0=\alpha\beta Q-(\gamma-(\alpha+\beta+1)x)\dfrac{\mathrm dQ}{\mathrm dx}\\&\qquad-(x-xx)\dfrac{\mathrm{dd}Q}{\mathrm dx^2}\\[4pt]
\mathrm{III.}\quad&0=\alpha\beta R-(\gamma-(\alpha+\beta+1)x)\dfrac{\mathrm dR}{\mathrm dx}\\&\qquad-(x-xx)\dfrac{\mathrm{dd}R}{\mathrm dx^2}
\end{aligned}
\]
Multiplying the first equation by \(Q\) and the second by \(P\), and subtracting, gives
\[
\begin{aligned}&0=(\gamma-(\alpha+\beta+1)x)\frac{Q\mathrm dP-P\mathrm dQ}{\mathrm dx}\\&+(x-xx)\frac{Q\mathrm{dd}P-P\mathrm{dd}Q}{\mathrm dx^2}\end{aligned}
\]
This equation, when multiplied by \(x^{\gamma-1}(1-x)^{\alpha+\beta-\gamma}\), becomes integrable and yields
\SourceLeaf{231}
\[
[94]\qquad A=x^\gamma(1-x)^{\alpha+\beta+1-\gamma}\frac{Q\mathrm dP-P\mathrm dQ}{\mathrm dx}
\]
Exactly the same procedure gives
\[
\begin{array}{rl}
[95]&B=x^\gamma(1-x)^{\alpha+\beta+1-\gamma}\dfrac{R\mathrm dQ-Q\mathrm dR}{\mathrm dx}\\[11pt]
[96]&C=x^\gamma(1-x)^{\alpha+\beta+1-\gamma}\dfrac{R\mathrm dP-P\mathrm dR}{\mathrm dx}
\end{array}
\]
The constants \(A,\ B,\ C\) are easily determined by the following method.

For \(x=0\), \(P=1\); further, \(x^\gamma Q=xF(\alpha+1-\gamma,\beta+1-\gamma,2-\gamma,x)\) becomes \(=0\) for \(x=0\); its differential divided by \(\mathrm dx\), namely \(\gamma x^{\gamma-1}Q+x^\gamma\frac{\mathrm dQ}{\mathrm dx}\), however, becomes \(=1\); hence \(\frac{x^\gamma\mathrm dQ}{\mathrm dx}=1-\gamma\) follows for \(x=0\), and therefore
\[
A=\gamma-1
\]
To determine \(B\) and \(C\) as well, let us return to the equation
\[
R=f(\alpha,\beta,\gamma)P+f(\alpha+1-\gamma,\beta+1-\gamma,2-\gamma)Q
\]
which, on differentiation, gives
\[
\frac{\mathrm dR}{\mathrm dx}=f(\alpha,\beta,\gamma)\frac{\mathrm dP}{\mathrm dx}+f(\alpha+1-\gamma,\beta+1-\gamma,2-\gamma)\frac{\mathrm dQ}{\mathrm dx}
\]
Multiplying the first by \(\frac{\mathrm dQ}{\mathrm dx}\) and the second by \(Q\), and subtracting, gives
\[
\frac{Q\mathrm dR-R\mathrm dQ}{\mathrm dx}=f(\alpha,\beta,\gamma)\frac{Q\mathrm dP-P\mathrm dQ}{\mathrm dx}\quad\mbox{therefore}
\]
\[
B=(1-\gamma)f(\alpha,\beta,\gamma)=\frac{\Pi(\alpha+\beta-\gamma)\Pi(1-\gamma)}{\Pi(\alpha-\gamma)\Pi(\beta-\gamma)}
\]
Similarly, multiplying the former equation by \(\frac{\mathrm dP}{\mathrm dx}\) and the latter by \(P\), subtraction gives
\[
\frac{R\mathrm dP-P\mathrm dR}{\mathrm dx}=f(\alpha+1-\gamma,\beta+1-\gamma,2-\gamma)\frac{Q\mathrm dP-P\mathrm dQ}{\mathrm dx}
\]
and therefore
\[
\begin{aligned}&C=(\gamma-1)f(\alpha+1-\gamma,\beta+1-\gamma,2-\gamma)\\&=\frac{\Pi(\alpha+\beta-\gamma)\Pi(\gamma-1)}{\Pi(\alpha-1)\Pi(\beta-1)}\end{aligned}
\]
If preferred, these three equations can also be written so that the derivatives \(\frac{\mathrm dP}{\mathrm dx},\ \frac{\mathrm dQ}{\mathrm dx},\ \frac{\mathrm dR}{\mathrm dx}\) are expressed by finite functions; thus, for example, formula 96 becomes
\SourceLeaf{232}
\[
\begin{aligned}
{}[97]\quad&\frac1\gamma F(\alpha,\beta,\alpha+\beta+1-\gamma,1-x)\\&\qquad F(\alpha+1,\beta+1,\gamma+1,x)\\
&+\frac1{\alpha+\beta+1-\gamma}F(\alpha+1,\beta+1,\alpha+\beta+2-\gamma,1-x)\\&\qquad F(\alpha,\beta,\gamma,x)\\
&=\frac{\Pi(\alpha+\beta-\gamma)\Pi(\gamma-1)}{\Pi\alpha\Pi\beta}x^{-\gamma}(1-x)^{\gamma-\alpha-\beta-1}
\end{aligned}
\]
\SixthArticle{52}
Denoting the function $F(-\alpha,-\beta,1-\gamma,x)$ by $S$ will give
\[
0=\alpha\beta S-(1-\gamma+(\alpha+\beta-1)x)\frac{\mathrm dS}{\mathrm dx}-(x-xx)\frac{\mathrm{dd}S}{\mathrm dx^2}
\]
Combining this equation with I of the preceding article gives
\[
0=\alpha\beta\left(\frac{S\mathrm dP+P\mathrm dS}{\mathrm dx}\right)
-(1-2x)\frac{\mathrm dP}{\mathrm dx}\mathbin{.}\frac{\mathrm dS}{\mathrm dx}
-(x-xx)\frac{\mathrm dS\mathrm{dd}P+\mathrm dP\mathrm{dd}S}{\mathrm dx^2}
\]
which is integrable and yields
\[
\mathrm{Const.}=\alpha\beta PS-(x-xx)\frac{\mathrm dP}{\mathrm dx}\mathbin{.}\frac{\mathrm dS}{\mathrm dx}
\]
The value of the constant quantity $=\alpha\beta$, as follows of itself from $x=0$. If you prefer a finite form, you have
\[
\begin{aligned}
{}[98]\quad&F(\alpha,\beta,\gamma,x)F(-\alpha,-\beta,1-\gamma,x)\\
&-\frac{\alpha\beta}{\gamma-\gamma\gamma}\\&\qquad(x-xx)F(\alpha+1,\beta+1,\gamma+1,x)\\&\qquad F(1-\alpha,1-\beta,2-\gamma,x)=1
\end{aligned}
\]
Transforming each of the four functions here according to formula 82, and then writing $\gamma-\alpha,\gamma-\beta$ for $\alpha,\beta$, you will have
\[
\begin{aligned}
{}[99]\quad&(1-x)F(\alpha,\beta,\gamma,x)F(1-\alpha,1-\beta,1-\gamma,x)\\
&-\frac{\gamma-\alpha\mathbin{.}\gamma-\beta}{\gamma-\gamma\gamma}\\&\qquad xF(\alpha,\beta,\gamma+1,x)F(1-\alpha,1-\beta,2-\gamma,x)=1
\end{aligned}
\]
